This section reports the failure mode that motivated this paper's central
fix, largely as originally measured (on the pre-fix system), because the
finding itself --- and the diagnostic process that led from a suspicious
number to a specific one-line root cause --- is the useful part.
\S\ref{sec:controlled} and \S\ref{sec:results} report the post-fix
measurement.

The dominant failure mode observed on the pre-fix system was
over-fragmentation of the knowledge graph driven by a high rate of
predicted CONTRADICTS edges, most of them spurious. On SMRITI-Reference,
Phase 6 classified 1{,}539 of 4{,}847 validated candidate pairs (31.8\%)
as CONTRADICTS, and Phase 7's partitioning collapsed 569 claim nodes into
472 partitions --- an 82.9\% singleton rate. Given the corpus was
deliberately constructed with exactly three genuinely contradictory
document pairs, a rate this high could not be explained by genuine
semantic conflict.

To isolate the mechanism, we drew a sample of 60 candidate pairs in which
one claim comes from a deliberately-planted adversarial document and the
other comes from an \emph{unrelated} document elsewhere in the corpus ---
pairs that should essentially always be NEUTRAL. SMRITI classified 11 of
these 60 pairs (18.3\%) as CONTRADICTS; both independent annotation passes
agreed that all 11 (100\%) were false positives.

\paragraph{Root cause.} Inspecting the resolver directly
(\texttt{retrieval/classification/resolver.py}) rather than only its
output distribution located the mechanism precisely: the CONTRADICTS rule
required \texttt{contradiction\_score > entailment\_score} by a margin,
but never compared against \texttt{neutral\_score} at all. On a candidate
pair drawn from embedding-retrieval rather than a curated, topically-
coherent test set, an off-the-shelf NLI model frequently assigns its
\emph{highest} probability to neutral (correctly recognizing the pair as
unrelated) while still assigning contradiction a higher score than
entailment (incorrectly, but not the model's dominant judgment) --- and the
old resolver committed to CONTRADICTS anyway, because it never looked at
the class the model actually favored. This is not a fundamental limitation
of the NLI model; it is an integration bug in how a three-class output was
reduced to a decision. \S\ref{sec:system} describes the fix (require the
winning class to be the true arg-max of all three scores, plus a
topical-relatedness gate specific to CONTRADICTS), and \S\ref{sec:controlled}
and \S\ref{sec:results} report that it holds on both a fully objective
corpus and the re-annotated SMRITI-Reference sample.

\paragraph{Downstream consequence for reliability scoring.} Because Phase
8's reliability signals include topology strength computed over each
claim's partition, the fragmentation above directly and mechanically
suppressed reliability scores system-wide: a claim isolated into its own
singleton partition has no corroborating neighbors by construction,
regardless of how well-established the claim itself is. At the time of
this measurement, Phase 8 fused topology strength directly into
reliability\_index; \S\ref{sec:system} and \S\ref{sec:limitations}
describe a later, separate rectification (P1-4) that reports graph
importance as its own index rather than blending it into reliability at
all, so this specific mechanical suppression path no longer applies to
reliability\_index post-split (importance\_index inherits it instead,
which is the correct scope for a graph-structure effect). This
fragmentation is not a Phase 8 calibration bug in isolation --- it is a
symptom propagating upward from Phase 6, and improves automatically as
Phase 6's precision improves; \S\ref{sec:limitations} reports the current
partitioner's own partition count and singleton rate on the post-fix
corpus (55.6\%, substantially lower than the 82.9\% above, though not
directly comparable since the current figure isolates partitioning
\emph{algorithm} choice on an already-precision-fixed edge set, not
residual Phase 6 CONTRADICTS precision).

\paragraph{Secondary failure mode: claim-extraction noise from document
metadata.} Phase 4 previously extracted a claim from every semantic
sentence Phase 3 produced, including purely bibliographic or structural
sentences. \S\ref{sec:system}'s assertion classifier addresses this
directly; \S\ref{sec:results} quantifies the resulting change in extracted
claim volume and validity.

\subsection{Error Taxonomy}
\label{sec:error-taxonomy}

Rather than propose failure categories in the abstract, we derive a
taxonomy directly from the 482 core row-level outcomes in
Table~\ref{tab:controlled} (SMRITI-Controlled-v2, the 713-pair,
15-category scale-up of the zero-annotation corpus, post-bidirectional-NLI-
rewrite) plus the pipeline-level failure modes already characterized
above on SMRITI-Reference. Every entry below is traceable to specific,
real pair IDs or a specific, reproducible category-level rate; we do not
report a category we did not actually observe. An earlier draft of this
taxonomy was built from a 32-pair pilot version of this corpus; several
of its findings (E2, E3, E7) turned out, once corpus size grew by roughly
20$\times$, to be either not reproduced or far more severe than the
pilot's small sample suggested. We report the revision itself as a
finding: a construction-time-gold corpus removes annotation noise, but
not sampling noise, and $n=1$--6 anecdotes in the pilot version of this
section should not have been read as rates.

\begin{itemize}
  \item \textbf{E1 -- Retrieval miss.} The true pair is never surfaced as
    an NLI candidate at all, so no classification is possible. On the
    sanitized corpus's frozen TEST split this is the sole determinant of
    each category's retrieval-recall column in Table~\ref{tab:controlled}.
    All six CONTRADICTS subtypes now reach 100\% retrieval recall on
    TEST (the corpus's one remaining CONTRADICTS TEST miss, in
    \emph{temporal}, is a classification miss given retrieval, not a
    retrieval miss --- see the per-class table). Retrieval miss remains
    real for \emph{SUPPORTS} (25.0\% miss rate on TEST, 3 of 12) and for
    \emph{EQUIVALENT} (28.6\%, 4 of 14). Root cause is Phase 5/6
    candidate generation (cosine-similarity thresholding), upstream of
    the resolver fix this paper centers on.
  \item \textbf{E3 -- RETRACTED: an earlier draft's claim that SUPPORTS
    classification collapses to zero correct labels on the sanitized
    corpus's frozen TEST split did not match the frozen results
    artifact.} The claim previously here read: ``0 of 9 retrieved TEST
    SUPPORTS pairs are correctly labeled SUPPORTS.''
    \texttt{evaluation/controlled/v2/results.json}'s per-pair rows show
    the opposite: all 9 of the 9 retrieved TEST SUPPORTS pairs are
    correctly labeled SUPPORTS (precision 90.0\%, recall 75.0\%, F1
    81.8\%, DEV: precision 85.0\%/recall 87.2\%/F1 86.1\%). SUPPORTS
    classification given retrieval is not a resolver weakness on this
    corpus; the residual TEST recall gap is retrieval (E1: 3 of 12
    SUPPORTS pairs never reach the classifier), not classification. We
    keep this entry, retracted rather than deleted, so a reader comparing
    this paper against an earlier draft or an earlier audit
    round can see exactly what changed and why (\S\ref{sec:controlled}
    finding 3 has the full per-pair reconciliation). SUPPORTS
    classification on \emph{real}, naturally-occurring prose remains
    weak (\S\ref{sec:limitations}) --- that is a separate finding, on a
    separate corpus, and is not affected by this retraction.
  \item \textbf{E7 -- Same-entity multi-attribute conflation is the
    dominant failure mode of the entire hard-NEUTRAL extension.} A
    genuinely NEUTRAL pair describing two \emph{compatible} attributes of
    the same named entity (e.g.\ ``the \{X\} uses Redis for response
    caching'' / ``uses an in-memory LRU cache for response caching'') is
    frequently misclassified CONTRADICTS. On the sanitized corpus's
    frozen TEST split (17 purpose-built multi-attribute pairs), all 17
    are retrieved (100\%, confirming the relatedness gate correctly
    recognizes these as topically related --- that is its entire job),
    and only 5 of 17 (29.4\%) are correctly labeled NEUTRAL; the other 12
    are all misclassified CONTRADICTS. (An earlier draft of this entry
    reported ``zero are correctly labeled NEUTRAL''; the frozen results
    artifact shows 5, not 0, and this entry is corrected accordingly ---
    the finding survives at a less extreme magnitude.) Neither the
    relatedness gate nor the NLI model reliably distinguishes ``two
    attributes describing the same entity'' from ``two mutually exclusive
    claims about the same entity''; ``X uses A'' and ``X uses B'' most
    often read as a factual disagreement about which is true, when both
    can be. This is this corpus's most consequential open item for
    future resolver work, and it survives corpus sanitation unchanged in
    kind --- it is a genuine resolver/gate gap, not a symptom of the
    grammar or exclusivity defects fixed elsewhere in this corpus.
  \item \textbf{E9 -- Implicit-exclusivity phrasing weakening the
    \emph{direct} and \emph{entity} CONTRADICTS subtypes.} \textbf{Fixed,
    not merely disclosed.} A second audit round identified that an
    earlier draft's \emph{direct} subtype (``runs on \{A\}'' vs.\ ``runs
    on \{B\}'') and the original pilot corpus's \emph{entity} subtype
    (``developed by Team Alpha'' vs.\ ``developed by Team Beta'') did not
    actually satisfy CONTRADICTS's own definition: a system can run on
    both architectures, and a project can have multiple contributing
    teams, so the two claims in each pair could both be true at once ---
    the corpus was asserting a contradiction the text did not itself
    establish. \S\ref{sec:controlled} describes the fix: every CONTRADICTS
    pair in the regenerated corpus now carries an explicit exclusivity
    premise (``primary deployment platform,'' ``developed \emph{solely}
    by'') and \texttt{contradiction\_basis}/\texttt{exclusivity\_type}
    metadata, checked by \texttt{validate\_corpus.py}. Post-fix,
    \emph{direct} and \emph{entity} both reach 100\% TEST retrieval
    recall with 100\% classification accuracy given retrieval --- we
    record this as a fixed defect, discovered by audit, rather
    than continuing to report the pre-fix numbers as an ongoing
    corpus-construction caveat.
  \item \textbf{E10 -- No measurable abstention on genuinely ambiguous
    input, and confidence/margin/entropy diagnostics show this is a
    benchmark-design gap, not a threshold-tuning one.} \textbf{Open,
    reclassified from a prior draft's ``substantially achieved'' framing
    (P0-E, audit round).} \texttt{ResolutionStatus.ABSTAINED}
    (\S\ref{sec:system}, audit round P0-4) is a formally defined,
    reachable resolver outcome, but on the frozen TEST split's 15 pairs
    purpose-built with hedged, non-committal language (``build time
    varies depending on project size,'' ``some large projects report
    longer build times'') so that no confident label is well-supported by
    the text, the resolver abstains zero times: all 15 are retrieved as
    candidates and confidently classified NEUTRAL. Confidence/margin/
    entropy diagnostics on these same 15 pairs' own recorded NLI scores
    (Table~\ref{tab:abstention-diagnostic}, \S\ref{sec:controlled}
    finding 9) show why no threshold adjustment would fix this: mean
    neutral-class confidence is 98.9\%/98.7\% (A$\to$B/B$\to$A), mean
    margin over the runner-up class is 97.5--97.8 points, and the
    \emph{lowest} neutral confidence observed across all 15 pairs, either
    direction, is still 92.5\%. Per the selective-prediction framing this
    same review recommended, we report this as \textbf{0\% abstention
    rate / 100\% coverage} on a population explicitly constructed to
    warrant abstention, not as ``UNKNOWN classification accuracy''
    against some single correct label these pairs do not have --- but we
    go further than that framing alone: \texttt{ResolutionStatus.ABSTAINED}
    is a \emph{confidence-based} mechanism (fires when the model is
    insufficiently confident to commit), not a \emph{semantic-ambiguity}
    detector (recognizing that a claim pair is genuinely hedged or
    underspecified, independent of model confidence), and this benchmark
    tests the latter while only the former is implemented. This is not
    evidence \texttt{ABSTAINED} is unreachable in general (SMRITI-Reference's
    own pipeline does emit it on other inputs, and unit tests correctly
    verify the low-confidence branch fires when confidence genuinely is
    low), but it is evidence that hedged natural language does not
    reliably produce low NLI confidence in this model, so a benchmark
    built from hedged language does not exercise abstention no matter how
    the threshold is tuned.
  \item \textbf{E11 -- More than half of all CONTRADICTS verdicts fire on
    candidate pairs the corpus never intended to test.} \textbf{Open, new
    finding (audit round P0-9).} Measured by retaining every
    NLI-scored candidate (\texttt{--env eval}), not only the ones
    surviving production's NEUTRAL/UNKNOWN filters: of 1{,}089 total
    candidates Phase 6 scored across the full corpus, only 496 (45.6\%)
    correspond to a manifest-intended pair at all, and of all
    CONTRADICTS verdicts, 383 (55.2\%) are on non-intended pairs. This is
    not the same failure as E4 (the original arg-max/relatedness-gate
    bug, which produced false CONTRADICTS between claims sharing no real
    topic): these non-intended pairs were retrieved because embedding
    similarity genuinely found them close, most plausibly because many
    template families reuse a limited pool of attribute/verb phrases
    across dozens of different entities, so two \emph{different} intended
    pairs' claims can end up lexically closer to \emph{each other} than
    either is to some unrelated claim. Whether these 383 calls are
    themselves correct is not established (they were never given a
    construction-time label), but their volume relative to the 119
    intended, scored CONTRADICTS trials on TEST means gold-pair-only precision
    (finding 1, \S\ref{sec:results}) is measuring a curated slice of a
    much larger and largely unaudited set of real verdicts the same run
    produces.
  \item \textbf{E4 -- Spurious cross-document contradiction} (pre-fix
    only; \textbf{fixed}, \S\ref{sec:bugs} Defect 5). Unrelated pairs
    resolved to CONTRADICTS because the old resolver never checked the
    NLI model's dominant class. Measured pre-fix at 18.3\% on a
    targeted 60-pair sample, 100\% of which both annotation passes
    confirmed false; measured post-fix at 0/6 on SMRITI-Controlled's
    NEUTRAL pairs (E4's absence is itself the finding in
    \S\ref{sec:controlled}, point (2)).
  \item \textbf{E5 -- Non-declarative claim noise} (pre-fix only;
    \textbf{fixed}, \S\ref{sec:system} assertion classifier). Headings,
    bibliographic lines, and list fragments constructed into claims with
    no assertion content, inflating claim volume and diluting every
    downstream signal that treats claim count as a denominator.
  \item \textbf{E6 -- Rhetorical stance conflict misread as neutral.}
    \textbf{Open, newly discovered post-fix.} On SMRITI-Reference's real
    prose, 22 of 25 deliberately-planted hard-contradiction pairs are
    predicted NEUTRAL, not because the resolver malfunctions but because
    the NLI model itself assigns neutral $>0.97$ probability to these
    pairs (\S\ref{sec:results-annotation} inspects the raw scores
    directly). The disagreement in these pairs is expressed as
    argumentative prose about a shared topic (``U-Net is an outdated
    architectural choice'') rather than direct propositional negation,
    which a model trained on SNLI/MultiNLI/ANLI-style propositional pairs
    is never asked, during training, to read as contradiction. Unlike E4,
    this is not an integration bug: the resolver is faithfully reporting
    the model's own dominant class. It requires a stance-detection-
    specific model, not a resolver change and not a different calibration
    of the same model --- calibration would only make an already-correct-
    for-this-task probability more trustworthy, not teach the model a
    task it was never trained on. Recorded as unresolved in
    \S\ref{sec:limitations}.
\end{itemize}

E1, E6, E7, E10, and E11 remain open (E1 is a retrieval-stage problem
with an uneven, category-dependent rate; E3 is RETRACTED, not open --- SUPPORTS
classification is strong given retrieval (\S\ref{sec:controlled} finding
3), and the residual TEST gap is a retrieval, not classification,
problem (E1); E6 is an NLI task mismatch between rhetorical
and propositional disagreement, not a calibration gap; E7 is this
evaluation's single largest source of classification error on the
hard-NEUTRAL extension, a same-entity multi-attribute conflation the
relatedness gate cannot catch by design; E10 is a confidence-based-
abstention-vs-semantic-ambiguity-detection gap --- 0\% abstention rate on
text engineered to warrant it, with confidence/margin/entropy
diagnostics confirming this is not a threshold-tuning fix away, not
``inaccuracy'' against a label that does not exist; and E11 is a newly
measured, large volume of CONTRADICTS verdicts on candidate pairs outside
this corpus's own intended test set (not a false-positive rate --- these
pairs were never given a construction-time label, so their correctness
is not established either way), invisible to any gold-pair-only
metric); E9 is a corpus-construction defect a second audit round
found and this revision fixes, not merely discloses; E4--E5 are the two
mechanisms this paper's central fixes target and close on the evidence
available here. We do not claim this
list is exhaustive over failure modes SMRITI could exhibit on other
corpora --- only that it is exhaustive over failure modes we actually
observed on the two corpora reported in this paper.
